\section{Fixed-phase powers and fractional repetitions}
\label{sec:fixed-phase}

The conditions on the phase in Section~\ref{sec:phase-counts}
differ from those in the following two results. The first gives a
small power saving at a prescribed phase for every quadratic
irrational. The second gives prescribed-phase counts for a
shortened comparison in arbitrary degree.

\begin{theorem}\label{thm:fixed-quad}
Fix a real quadratic irrational \(\alpha\) and \(s\ge2\).
There is \(\eta>0\) such that, with \(A=B^{s-\eta}\), the
following assertion holds for fixed bounded \(I\) of positive
length. Let \(J_h\subset[0,1)\) be a half-open interval of length
at least \(B^{-\eta}\) and \(1\le q\le B^\eta\). For all real
\(\theta\) and
\[
 \beta=(u b^m\alpha+v)/d,\quad
 1\le |u|,d\le B^\eta,\quad 0\le m\le\eta h,
 \qquad u,v,d,m\in\Z,
\]
the number of integers \(n/A\in I\) for which \(\{\theta+n\beta\}\)
starts with \(V^s\), \([V]_b/B\in J_h\), and
\([V]_b\equiv w_0\bmod q\), is
\[
 (1+o(1))\, \frac{AB |I|\,|J_h|}{M q}.
\]
The error is uniform in the stated ranges. In particular at phase
zero there are powers of least period \(h\) with positive multipliers
\(a\equiv1\bmod b,\ a=o(B^{s-\eta})\). For squares, the root can be
prescribed to begin with \(U\) or \(cU\) and to end with
\(\overleftarrow U\), where \(U\) is any word of length
\(\lfloor\eta h\rfloor-1\). These choices give an even or an odd
palindrome at the seam, respectively.
\end{theorem}

\begin{proof}
We first show there exists \(\zeta>0\) such that
\[
 0<|k|\le B^{s+\zeta},\quad |r|\le B^\zeta,\quad
 k\equiv r\bmod D,\quad \|k\alpha\|\le B^{-s+\zeta}
\]
has no solution for large \(h\). Otherwise take a sequence of
counterexamples with \(\zeta\to0\), \(h\to\infty\), and \(k>0\).
The quadratic estimate \(\|k\alpha\|\gg 1/k\) gives
\(\log k=(s+o(1))\log B\).
Let \(a_0 X^2+a_1X+a_2\) be an integer minimal polynomial,
\(\Delta=a_1^2-4a_0a_2\), and \(2a_0\alpha+a_1=\sigma\sqrt\Delta\)
with \(\sigma=\pm1\). For the nearest integer \(P\) to \(k\alpha\),
put \(Z=2a_0P+a_1 k+\sigma k\sqrt\Delta\). Its conjugate is
\(\overline{Z}=2a_0(P-k\alpha)\), and
\(Z-\overline Z=2\sigma k\sqrt\Delta\).
Hence \(|Z|\asymp k\) and \(0<|Z\overline{Z}|\ll B^{2\zeta}\).
Use a fixed norm-one Pell unit \(\epsilon>1\) of
\(\Z[\sqrt\Delta]\) and write \(Z=z\epsilon^j\) with
\(1\le |z|<\epsilon\).
Then \(j\log\epsilon=(s+o(1))\log B\), and the coefficients \(z\),
\(\overline z\), and \(r\) have logarithmic height \(o(\log B)\).

Take \(S\) to consist of the archimedean places and those over
primes dividing \(b\) in \(\Q(\sqrt\Delta)\).
The two arguments \(X=B,\ Y=\epsilon^j\) are \(S\)-units of
maximum logarithmic height \(H\asymp\log B\).
The moving polynomials \(X-1\) and
\(zY^2-2\sigma\sqrt\Delta\,rY-\overline z\) are coprime of
fixed positive degrees with nonzero constant terms and height
\(o(H)\). The second value equals
\(2\sigma\sqrt\Delta\,Y(k-r)\), divisible by \(B-1\) at the
finite places in the ring of integers. Their logarithmic gcd
is therefore at least \(\log(B-1)\).
Grieve--Wang's moving gcd theorem
\cite[Theorem~1.2]{GrieveWang}, with sufficiently small
gcd tolerance, implies on an infinite subsequence that
\(B^a(\epsilon^j)^d\) has height \(o(H)\) for a fixed nonzero
integer pair \((a,d)\).
If \(a\ne0\), a valuation over a prime dividing \(b\) gives
height \(\gg h\), as the Pell unit has valuation zero.
If \(a=0\), the height is \(|d|j\operatorname{ht}(\epsilon)\gg h\).
Both cases contradict the conclusion. This proves the exclusion.

Choose \(0<\eta<\min(\zeta/32,1/16)\).
Let \(W=\{w:w/B\in J_h,\ 0\le w<B,\ w\equiv w_0\bmod q\}\).
It has cardinality \(B|J_h|/q+O(1)\gg B^{1-2\eta}\).
Smooth each digit cylinder at \(Lw/M\) on scale \(1/M\)
using a single compactly supported \(C^\infty\) function
\(\phi(Mx-Lw)\), periodized in \(x\). Their sum has coefficient
\[
 p\,\widehat\phi(k/M)S(k),\quad
 p=|W|/M,\qquad
 S(k)=|W|^{-1}\sum_{w\in W}e(-kLw/M).
\]
For the multiplier smoothing \(g(n/A)\), put
\(G_A(t)=\sum_n g(n/A)e(nt)\). Its magnitude is bounded by
\(O(A(1+A\|t\|)^{-8})\).
For \(1\le |k|\le T_1\) the quadratic estimate for \(\alpha\)
makes the sample points \(k\beta\bmod1\) separated by
\(\gg 1/(T_1 B^{3\eta})\).
Counting points in bands of width \(1/A\) gives, for \(T_1\ge A\),
\[
 \begin{aligned}
 \sum_{|k|\le T_1}|G_A(k\beta)|&\ll B^{3\eta}T_1,\\
 \sum_{\substack{|k|\le T_1\\ A\|k\beta\|>Z_1}}
       |G_A(k\beta)|&\ll B^{3\eta}T_1(1+Z_1)^{-7}
                  \qquad (Z_1\ge2).
 \end{aligned}
\]
For \(0<|k|\le K=MB^\eta\) with \(|S(k)|\ge B^{-8\eta}\),
the geometric sum gives
\(\|qkL/M\|\ll B^{-1+10\eta}\).
Taking the nearest integer \(l\) gives
\(r_1=qk-lD\) of magnitude \(O(B^{10\eta})\), using \(M=LD+1\).
If \(\|k\beta\|\le B^{-s+\zeta-2\eta}\), then the frequency
\(qk u b^m\) has magnitude at most \(B^{s+4\eta}\), residue
\(r_1 u b^m\) of magnitude \(O(B^{12\eta})\),
and its distance after multiplication by \(\alpha\) is
at most \(qd\|k\beta\|\le B^{-s+\zeta}\).
This is excluded by the first paragraph. Thus
\(A\|k\beta\|>B^{\zeta-3\eta}\) for these \(k\).
Relative to \(pA\), the error for \(|k|>K\) is
\(O(B^{3\eta}(M/A)(K/M)^{-7})=O(B^{-3\eta})\) by dyadic
summation and decay of \(\widehat\phi\).
For \(|k|\le K\) the small-\(|S(k)|\) and large-\(|S(k)|\)
errors are bounded by \(O(B^{-3\eta})\) and
\(O(B^{5\eta-7(\zeta-3\eta)})\), respectively.
The zero term yields \(pA\int g\int\phi+o(pA)\).
Sandwiching the half-open interval and the cylinders with
smooth functions proves the asserted exact count.

For the special multipliers take \(\beta=b\alpha,\theta=\alpha\).
Distinct tested multipliers have circle separation \(\gg 1/A>1/M\)
for large \(h\). There are \(\gg B^{1-3\eta}\) distinct roots
even with the prescribed templates, because the prefix interval
has width at least \(B^{-\eta}\) and the suffix prescribes one
residue with modulus at most \(B^\eta\).
Discard \(O(B^{1/2})\) roots with a period at most \(h/2\) to
obtain least period \(h\) for the full power. Allowing any fixed
positive proportion bound on \(a/A\) and taking a minimum yields
the little-oh assertion.
\end{proof}

\begin{theorem}\label{thm:fractional-fixed}
Fix any real algebraic irrational \(\alpha\), \(0<\xi<1\),
\(\delta,C>0\), and put \(r=\lfloor\xi h\rfloor\), \(E=b^r\),
\(A=E^2B^\delta\). For fixed bounded \(I\) of positive length,
among integers \(n/A\in I\) the number for which the first \(r\)
digits of \(\{\theta+n b^m\alpha\}\) agree with its length-\(r\)
block after \(h\) places is \((1+o(1))|I|A/E\), uniformly for
\(\theta\in\R\) and \(0\le m\le Ch\).
If \(0<\delta<1-\xi\) and
\(\xi+\delta>\max(1/2,1-\xi)\), a positive proportion of the hits
have least period \(h\) on the entire length-\((h+r)\) segment.
\end{theorem}

\begin{proof}
Let \(\kappa=\delta/16\), \(\lambda=\delta/4\).
For a fixed finite set of integers \(j\), all nonzero
\(k=jE+lD\), \(|l|\le EB^\kappa\), satisfy
\(\|k b^m\alpha\|>A^{-1}B^\lambda\) for large \(h\).
For \(l=0\) this is \eqref{eq:phase-ridout}, applied to
\(j b^{m+r}\). For \(l\ne0\) use the integer vector
\[
     (X_1,X_2,X_3,P)=(b^{m+h}l,b^m l,b^mE,P)
\]
where \(P\) is the nearest integer to \(k b^m\alpha\).
The real forms are the first three coordinates and
\(\alpha(X_1-X_2+jX_3)-P\). At primes dividing \(b\) use coordinate
forms. For a counterexample their product is at most
\(l^2 E^{-2}B^{-\delta+\lambda}\le B^{2\kappa-\delta+\lambda}\)
by cancellation of the base powers, and the height is \(B^{O(1)}\).
The \(S\)-adic Subspace Theorem \cite{Bilu} therefore puts an
infinite subsequence of counterexamples with fixed \(j\)
in a proper rational hyperplane. Divide its linear relation
by \(X_1\). The limit of \(P/X_1\) is \(\alpha\), while
\(X_2/X_1\) and \(X_3/X_1\) tend to zero. Irrationality
forces the coefficients of \(X_1,P\) to vanish.
The remaining relation gives \(l=c_0 E\) for a fixed nonzero
rational \(c_0\). Now \(k b^m\) equals
\(c_0 b^{m+r+h}+(j-c_0)b^{m+r}\).
Apply the same theorem on the vector
\((b^{m+r+h},b^{m+r},P)\) with error form
\(\alpha(c_0 X_1+(j-c_0)X_2)-P\).
The product is at most the error itself.
Dividing the new rational relation by its first coordinate
makes the limiting ratio of \(P\) equal to \(c_0\alpha\).
It forces all coefficients to vanish, a contradiction.

The exact block condition for \(x\in[0,1)\) is
\[
 -u\le t<1-u,\qquad u=\{Ex\},\quad t=E\langle Dx\rangle .
\]
Indeed \(\{Bx\}\) belongs to the same length-\(1/E\) cylinder
as \(x\) exactly when its signed circle displacement meets
these bounds. Use a trigonometric polynomial in \(u\), a smooth
test in \(t\), and a smooth multiplier test. The frequencies
are \(jE+lD\) and the nonzero terms with \(|l|\le EB^\kappa\)
have weighted error \(O(AB^{-N\lambda})\) by finite differences.
The tail is \(O(AB^{-\kappa(N-1)})\).
The zero term is \(A/E\) times the product of the test integrals
up to \(o(A/E)\). Taking \(N\) large and approximating the
region by rectangles as in Theorem~\ref{thm:selectable-counts}
yields the exact count, because its area is \(1\).

In the stated smaller range of \(\delta\),
\eqref{eq:phase-ridout} makes the hits use distinct
length-\((h+r)\) cylinders. For \(t_0<h\) a second period,
view a word of period \(t_0\) as labels modulo \(t_0\).
The equality constraints for period \(h\) join \(i\) to \(i+h\)
modulo \(t_0\) for \(0\le i<r\).
If \(r\ge t_0\), the graph has \(g=\gcd(h,t_0)\) components.
If \(r<t_0\), exactly \(\max(0,r-t_0+g)\) of the original shift
cycles remain complete, since their largest representatives are
\(t_0-g,\ldots,t_0-1\). There are then
\(t_0-r+\max(0,r-t_0+g)=\max(t_0-r,g)\) components.
Hence at most \(h b^{\max(h/2,h-1-r)}=o(A/E)\)
of all candidate words have a second period \(t_0<h\).
Discarding them proves the least-period assertion.
\end{proof}

Taking \(\xi=1-1/j\), \(\delta=1/(2j)\) and \(C=j\), and passing
through sufficiently large thresholds as \(j\to\infty\), gives
a comparison length \(r=h-o(h)\) with \(h-r\to\infty\).
The chosen positive multiplier is \(o(b^{h+r})\), and one
may take phase zero or a phase with \(m/h\to\infty\).
The comparison still leaves an intervening segment of \(h-r\)
digits.

